\section{Research intent and claim to test}
The primary objective is to improve detection of attacks for which only a small number of labelled examples are available. Improving every class or maximising aggregate macro-F1 is not the central claim. Under a limited context budget, the research question is which common and attack-specific examples help discriminate these attacks from benign traffic and other attacks while preserving existing detection performance.

TabPFN is a candidate because it conditions predictions on labelled context examples using a pretrained backbone. This permits context changes without updating backbone weights. It does not establish better few-label performance, lower total cost, or immediate online adaptation. XGBoost remains a strong baseline and, when model updates are studied, can use continued training as well as refitting. Preprocessing, context construction, caching, inference, calibration and any learned context or embedding must be charged to the relevant method.

\paragraph{Working hypothesis.} A context composition strategy can make scarce attack labels more useful than an unstructured context with the same information and memory budget. The resulting class-level benefit must be measured together with benign false alarms, other-class performance and computational cost. Scarcity is defined from available training labels, before examining the final evaluation outcomes.

\paragraph{Candidate mechanisms.} Diverse benign examples may better represent normal behaviour. Tail synthesis may broaden scarce attack coverage, with real-example duplication as a necessary control. Learned representations may improve class separation. Context optimisation may allocate limited context capacity more effectively. These are separate hypotheses, not established findings or a commitment to combine every mechanism. Relevant starting points include \href{https://proceedings.iclr.cc/paper_files/paper/2024/file/5d54d2df6ec8f7b920aa0fec9a6d1b2e-Paper-Conference.pdf}{LITO (ICLR 2024)}, \href{https://github.com/mlopezm/Supervised-contrastive-learning-over-prototype-label-embeddings}{supervised contrastive prototype embeddings}, and \href{https://arxiv.org/abs/2402.06971}{in-context data distillation}. The latter is associated with the ICLR 2024 ME-FoMo workshop, not the main conference.
